% Plantilla de Pandoc para los trabajos del Máster en Computación Cuántica (UNIR).
%
% Reproduce material/Plantilla_LATEX_FCC_MUCC_UNIR/plantilla.tex y carga su estilo
% estilo_unir-1.sty sin modificarlo (REGLAS.md, sección 8). Solo se añaden:
%   - las definiciones que necesita el LaTeX que genera Pandoc (listas, tablas,
%     figuras, citas de citeproc);
%   - tres ajustes técnicos que eliminan advertencias sin cambiar el diseño del modelo
%     (marcados con «ajuste técnico»).
\documentclass[11pt,a4paper,spanish]{book}
% Ajuste técnico: codificación T1. El modelo usa OT1, que no tiene versalita en negrita
% (babel la sustituye con una advertencia) y no separa bien las palabras con tilde.
\usepackage[T1]{fontenc}
\usepackage{estilo_unir-1}

% Ajuste técnico: la cabecera del modelo tiene dos líneas (\rhead{\theauthor\\}) y
% fancyhdr avisa de que 14pt no bastan; es la altura que fancyhdr aplicaría de todos modos.
\setlength{\headheight}{27.2pt}
% Ajuste técnico: la portada y el cuerpo empiezan ambos en la página 1; sin esto,
% hyperref crea dos destinos con el mismo nombre.
\hypersetup{hypertexnames=false}
% Ajuste técnico: la clase book iguala el final de todas las páginas (\flushbottom); con
% interlineado 1,5 y párrafos largos, LaTeX no siempre puede estirar la página y avisa
% (Underfull \vbox). Con \raggedbottom cada página termina en su última línea.
\raggedbottom

% --- Definiciones para el LaTeX de Pandoc -----------------------------------------
\providecommand{\tightlist}{%
  \setlength{\itemsep}{0pt}\setlength{\parskip}{0pt}}
\setlength{\emergencystretch}{3em}
% Citas y referencias generadas por citeproc (APA 7): lista con sangría francesa.
\NewDocumentCommand\citeproctext{}{}
\NewDocumentCommand\citeproc{mm}{%
  \begingroup\def\citeproctext{#2}\cite{#1}\endgroup}
\newlength{\cslhangindent}
\setlength{\cslhangindent}{1.27cm}
\newlength{\csllabelwidth}
\setlength{\csllabelwidth}{3em}
\newenvironment{CSLReferences}[2]
 {\begin{list}{}{%
  \setlength{\itemindent}{0pt}
  \setlength{\leftmargin}{0pt}
  \setlength{\parsep}{0pt}
  \ifodd #1
   \setlength{\leftmargin}{\cslhangindent}
   \setlength{\itemindent}{-1\cslhangindent}
  \fi
  \setlength{\itemsep}{#2\baselineskip}}%
  % APA 7: referencias alineadas a la izquierda (sin justificar); así las URL/DOI
  % largas no estiran la línea.
  \raggedright
  % El modelo carga natbib (author-year), cuyo \bibitem no acepta las entradas de
  % citeproc; dentro de esta lista cada entrada es un \item normal.
  \renewcommand{\bibitem}[2][]{\item}}
 {\end{list}}
\newcommand{\CSLBlock}[1]{\hfill\break\parbox[t]{\linewidth}{\strut\ignorespaces#1\strut}}
\newcommand{\CSLLeftMargin}[1]{\parbox[t]{\csllabelwidth}{\strut#1\strut}}
\newcommand{\CSLRightInline}[1]{\parbox[t]{\linewidth - \csllabelwidth}{\strut#1\strut}}
\newcommand{\CSLIndent}[1]{\hspace{\cslhangindent}#1}

%---------------------------
%título del trabajo y autor
%---------------------------
\title{Ejemplo de trabajo con el modelo UNIR}
\author{Julián Calderón Almendros}
\date{25 de septiembre de 2026}
\profesor{Rodrigo Gil-Merino y Rubio}

\begin{document}
\renewcommand{\listfigurename}{Índice de figuras}
\renewcommand{\listtablename}{Índice de tablas}
\renewcommand{\contentsname}{Índice de contenidos}
\renewcommand{\figurename}{Figura}
\renewcommand{\tablename}{Tabla}

\maketitle

\frontmatter
\tableofcontents

\chapter{Resumen}
Documento de ejemplo que muestra cómo se escribe un trabajo del máster
en Markdown con el modelo LaTeX de UNIR y referencias en formato APA 7.
Incluye citas narrativas y entre paréntesis, una fórmula display y un
apéndice. Sirve de referencia para los trabajos reales.

\vspace{\baselineskip}
{\bf Palabras clave:} modelo UNIR, formato APA, Markdown, LaTeX

\chapter{Abstract}
Sample document showing how a master's assignment is written in Markdown
with the UNIR LaTeX template and APA 7 references. It includes narrative
and parenthetical citations, a display formula and an appendix. It
serves as a reference for real assignments.

\vspace{\baselineskip}
{\bf Keywords:} UNIR template, APA style, Markdown, LaTeX

\mainmatter
\chapter{Introducción}\label{introducciuxf3n}

\section{Motivación y
justificación}\label{motivaciuxf3n-y-justificaciuxf3n}

La mecánica ondulatoria de Schrödinger (1926) describe el estado de un
sistema mediante una función de onda que obedece una ecuación de
autovalores. El debate sobre la completitud de esa descripción (Einstein
et~al., 1935) está en el origen de la información cuántica moderna
(Nielsen y Chuang, 2010).

\section{Planteamiento del trabajo}\label{planteamiento-del-trabajo}

La ecuación de Schrödinger independiente del tiempo para una partícula
de masa \(m\) en un potencial \(V(x)\) es

\[
\begin{aligned}
-\dfrac{\hbar^2}{2 \cdot m} \cdot \dfrac{d^2 \psi}{dx^2} + V(x) \cdot \psi(x) = E \cdot \psi(x).
\end{aligned}
\]

\section{Estructura del trabajo}\label{estructura-del-trabajo}

El capítulo siguiente resume el contexto; después se fijan los
objetivos, se desarrolla el trabajo y se exponen las conclusiones.

\chapter{Contexto y estado de la
cuestión}\label{contexto-y-estado-de-la-cuestiuxf3n}

El documento se compone con \LaTeX{} (Lamport, 1994).

\chapter{Objetivos}\label{objetivos}

\begin{enumerate}
\def\labelenumi{\arabic{enumi}.}
\tightlist
\item
  Mostrar la estructura del modelo.
\item
  Comprobar que las citas y la bibliografía siguen APA 7.
\end{enumerate}

\chapter{Desarrollo del trabajo}\label{desarrollo-del-trabajo}

Cuando se cita varias veces la misma obra con tres o más autores, APA
abrevia con «et al.» desde la primera cita (Einstein et~al., 1935, p.
778).

\chapter{Conclusiones}\label{conclusiones}

La estructura del modelo y el formato APA se generan desde una única
fuente en Markdown.

\chapter*{Bibliografía}\label{bibliografuxeda}
\addcontentsline{toc}{chapter}{Bibliografía}

\protect\phantomsection\label{refs}
\begin{CSLReferences}{1}{0}
\bibitem[\citeproctext]{ref-epr1935}
Einstein, A., Podolsky, B. y Rosen, N. (1935). Can quantum-mechanical
description of physical reality be considered complete? \emph{Physical
Review}, \emph{47}(10), 777--780.
\url{https://doi.org/10.1103/PhysRev.47.777}

\bibitem[\citeproctext]{ref-lamport1994}
Lamport, L. (1994). \emph{{LaTeX}: A document preparation system} (2.ª
ed.). Addison-Wesley.

\bibitem[\citeproctext]{ref-nielsen2010}
Nielsen, M. A. y Chuang, I. L. (2010). \emph{Quantum computation and
quantum information} (10.ª ed.). Cambridge University Press.
\url{https://doi.org/10.1017/CBO9780511976667}

\bibitem[\citeproctext]{ref-schrodinger1926}
Schrödinger, E. (1926). Quantisierung als Eigenwertproblem.
\emph{Annalen der Physik}, \emph{384}(4), 361--376.
\url{https://doi.org/10.1002/andp.19263840404}

\end{CSLReferences}

\appendix

\chapter{Apéndices}\label{apuxe9ndices}

Desarrollos auxiliares que no encajan en el cuerpo del trabajo.

\end{document}
